\documentclass[11pt,a4paper]{article}

\usepackage[margin=1in]{geometry}
\usepackage{booktabs}
\usepackage{amsmath}
\usepackage{amssymb}
\usepackage{graphicx}
\usepackage{hyperref}
\usepackage{xcolor}
\usepackage{caption}
\usepackage{subcaption}
\usepackage{float}

\hypersetup{
  colorlinks=true,
  linkcolor=black,
  citecolor=black,
  urlcolor=blue
}

\title{Reproduction and Extension of Half-Hop:\\
Large-Scale Message Passing via Graph Upsampling}
\author{Tejinderpal Singh}
\date{\today}

\begin{document}
\maketitle

\begin{abstract}
We reproduce key supervised experiments from Azabou et al.\ (ICML 2023) on
\textit{Half-Hop}, a graph upsampling augmentation that inserts intermediate
``slow nodes'' along edges to slow message passing.
We reimplement the Half-Hop operator, plug it into GCN, GraphSAGE, and GAT,
and report results on six heterophilic benchmarks (paper Table~2 style),
plus connectivity and initialization ablations on Texas.
As an extension beyond the paper, we evaluate Half-Hop on
\textbf{ModelNet10}, where 3D CAD meshes are converted to graphs and classified
with GCN/GraphSAGE with and without Half-Hop.
On heterophilic node classification, Half-Hop often improves accuracy.
On ModelNet10 mesh graphs, changes are small (about $+1\%$ for GCN and
$-0.5\%$ for GraphSAGE under a single-seed protocol).
\end{abstract}

\section{Introduction}
Message-passing neural networks can oversmooth node representations and
struggle when adjacent nodes have dissimilar labels (heterophily).
Half-Hop~\cite{azabou2023halfhop} addresses this by \emph{upsampling} the
input graph: for a directed edge $v_i \rightarrow v_j$, a slow node $\nu_k$
is inserted so that
\[
v_i \rightarrow \nu_k \rightleftharpoons v_j,
\]
with features
\[
\mathbf{x}_k = \alpha \mathbf{x}_i + (1-\alpha)\mathbf{x}_j.
\]
After message passing, slow nodes are discarded via a Boolean mask; only
original-node embeddings are used for prediction.
The method is model-agnostic: the same transform applies to GCN, GraphSAGE,
and GAT in supervised learning, and to view construction in self-supervised
settings.

Our goals are:
\begin{enumerate}
  \item implement Half-Hop faithfully and reproduce heterophilic supervised
        experiments;
  \item verify ablations on connectivity and slow-node initialization;
  \item test Half-Hop on \textbf{ModelNet10} 3D mesh graphs as an out-of-paper
        setting.
\end{enumerate}

\section{Method}
\subsection{Half-Hop transform}
Self-loops are never half-hopped.
For probability $p<1$, Half-Hop uses \emph{node-level} sampling: with
probability $p$, all incoming edges of a target node are half-hopped
(matching the official implementation).
We support the paper's connectivity variants (\texttt{proposed}, HH(1), HH(2))
and feature initializations (\texttt{linear}, \texttt{zero}, \texttt{random}).

\subsection{Models}
Baselines: GCN, GraphSAGE, GAT.
Half-Hop variants: HH-GCN, HH-GraphSAGE, HH-GAT, each applying Half-Hop before
the backbone and removing slow nodes after message passing.

\subsection{ModelNet10 extension}
ModelNet10 contains CAD meshes from 10 object categories.
We convert each mesh to a graph with PyTorch Geometric
\texttt{FaceToEdge}, use vertex coordinates as features, encode with
GCN or GraphSAGE (optionally after Half-Hop), then apply global mean pooling
and a linear classifier.
This is \emph{graph} classification (one label per mesh), unlike the paper's
node-classification tables.

\section{Experimental setup}
\paragraph{Heterophilic benchmarks.}
Texas, Wisconsin, Actor, Squirrel, Chameleon, Cornell
(Pei et al.\ splits; 10 train/val/test masks).
Hyperparameters for HH models follow paper Table~5 configs when available.
We report mean $\pm$ standard deviation of test accuracy at the best
validation epoch.

\paragraph{Ablations (Texas).}
Connectivity: proposed vs.\ HH(1) vs.\ HH(2).
Initialization: linear vs.\ zero vs.\ random.

\paragraph{ModelNet10.}
Train/test split from ModelNet; 10\% of training graphs held out for
validation.
Half-Hop uses $p=0.5$, $\alpha=0.5$; seed $0$; 100 epochs.
Hardware: Kaggle $2\times$ NVIDIA Tesla T4 (one model per GPU).

\section{Results}

\subsection{Heterophilic node classification}
Table~\ref{tab:hetero} summarizes our reproduction in the style of paper
Table~2.
Half-Hop improves GCN on all six datasets and improves GAT on five of six.
GraphSAGE gains are smaller; Cornell shows a decrease under our setup.

\begin{table}[H]
\centering
\caption{Heterophilic node classification (test accuracy \%, mean $\pm$ std
over 10 splits). $\Delta$ is HH $-$ baseline.}
\label{tab:hetero}
\resizebox{\textwidth}{!}{%
\begin{tabular}{lcccccc}
\toprule
Model & Texas & Wisconsin & Actor & Squirrel & Chameleon & Cornell \\
\midrule
Hom.\ level & 0.11 & 0.21 & 0.22 & 0.22 & 0.23 & 0.30 \\
\midrule
GCN & $58.65{\pm}3.61$ & $52.35{\pm}5.31$ & $28.68{\pm}1.02$ &
      $26.83{\pm}2.17$ & $37.83{\pm}2.25$ & $41.35{\pm}6.87$ \\
HH-GCN & $73.51{\pm}6.72$ & $55.49{\pm}5.92$ & $30.99{\pm}1.08$ &
         $29.63{\pm}1.07$ & $44.41{\pm}2.51$ & $62.97{\pm}8.74$ \\
$\Delta$ GCN & $+14.86$ & $+3.14$ & $+2.31$ & $+2.80$ & $+6.58$ & $+21.62$ \\
\midrule
GAT & $58.92{\pm}5.38$ & $50.20{\pm}5.33$ & $28.32{\pm}0.78$ &
      $29.21{\pm}1.39$ & $44.54{\pm}2.25$ & $43.24{\pm}8.26$ \\
HH-GAT & $70.81{\pm}9.08$ & $69.02{\pm}6.72$ & $28.49{\pm}1.85$ &
         $29.82{\pm}1.80$ & $44.41{\pm}2.53$ & $69.19{\pm}5.28$ \\
$\Delta$ GAT & $+11.89$ & $+18.82$ & $+0.17$ & $+0.61$ & $-0.13$ & $+25.95$ \\
\midrule
GraphSAGE & $76.76{\pm}6.14$ & $75.10{\pm}4.24$ & $33.88{\pm}1.01$ &
            $36.27{\pm}1.28$ & $50.48{\pm}2.51$ & $72.43{\pm}3.99$ \\
HH-GraphSAGE & $81.35{\pm}5.47$ & $79.22{\pm}7.63$ & $34.89{\pm}1.10$ &
               $37.10{\pm}1.32$ & $50.88{\pm}1.69$ & $67.84{\pm}6.43$ \\
$\Delta$ SAGE & $+4.59$ & $+4.12$ & $+1.01$ & $+0.83$ & $+0.40$ & $-4.59$ \\
\bottomrule
\end{tabular}%
}
\end{table}

Table~\ref{tab:vs-paper} compares our HH models to the paper's reported HH
numbers.
Some absolute gaps remain (especially Wisconsin/Squirrel/Chameleon),
consistent with baseline hyperparameter differences: the paper grid-searches
baselines, while we use published HH configs and simpler baseline defaults.

\begin{table}[H]
\centering
\caption{HH models vs.\ paper Table~2 (test accuracy \%).}
\label{tab:vs-paper}
\resizebox{\textwidth}{!}{%
\begin{tabular}{lcccccc}
\toprule
Dataset & Ours HH-GCN & Paper HH-GCN & Ours HH-GAT & Paper HH-GAT &
          Ours HH-SAGE & Paper HH-SAGE \\
\midrule
Texas & $73.51{\pm}6.72$ & $71.89{\pm}3.46$ &
        $70.81{\pm}9.08$ & $80.54{\pm}4.80$ &
        $81.35{\pm}5.47$ & $85.95{\pm}6.42$ \\
Wisconsin & $55.49{\pm}5.92$ & $79.80{\pm}4.30$ &
            $69.02{\pm}6.72$ & $83.53{\pm}3.84$ &
            $79.22{\pm}7.63$ & $85.88{\pm}3.99$ \\
Actor & $30.99{\pm}1.08$ & $35.12{\pm}1.06$ &
        $28.49{\pm}1.85$ & $36.70{\pm}0.92$ &
        $34.89{\pm}1.10$ & $36.82{\pm}0.77$ \\
Squirrel & $29.63{\pm}1.07$ & $47.19{\pm}1.21$ &
           $29.82{\pm}1.80$ & $46.35{\pm}1.86$ &
           $37.10{\pm}1.32$ & $45.25{\pm}1.52$ \\
Chameleon & $44.41{\pm}2.51$ & $60.24{\pm}1.93$ &
            $44.41{\pm}2.53$ & $61.12{\pm}1.83$ &
            $50.88{\pm}1.69$ & $62.98{\pm}3.35$ \\
Cornell & $62.97{\pm}8.74$ & $63.24{\pm}5.43$ &
          $69.19{\pm}5.28$ & $72.70{\pm}4.26$ &
          $67.84{\pm}6.43$ & $74.60{\pm}6.06$ \\
\bottomrule
\end{tabular}%
}
\end{table}

\subsection{Ablations on Texas}
Tables~\ref{tab:conn}--\ref{tab:init} match the paper's qualitative claims:
the proposed connectivity outperforms HH(1) and HH(2), and linear
interpolation outperforms zero and random slow-node features.

\begin{table}[H]
\centering
\caption{Connectivity ablation (Texas, HH-GCN).}
\label{tab:conn}
\begin{tabular}{lc}
\toprule
Scheme & Accuracy (\%) \\
\midrule
Proposed (HH) & $74.05{\pm}5.44$ \\
HH(1) & $70.00{\pm}6.80$ \\
HH(2) & $60.54{\pm}6.40$ \\
\bottomrule
\end{tabular}
\end{table}

\begin{table}[H]
\centering
\caption{Slow-node initialization ablation (Texas, HH-GCN).}
\label{tab:init}
\begin{tabular}{lc}
\toprule
Initialization & Accuracy (\%) \\
\midrule
Linear & $75.14{\pm}6.60$ \\
Zero & $68.92{\pm}5.73$ \\
Random & $66.22{\pm}6.89$ \\
\bottomrule
\end{tabular}
\end{table}

\subsection{Extension: ModelNet10 3D mesh classification}
Table~\ref{tab:modelnet} reports single-seed test accuracy at the best
validation epoch.
Half-Hop yields a small gain for GCN ($+0.99$) and a small drop for
GraphSAGE ($-0.55$).
Unlike heterophilic citation/web graphs, ModelNet meshes are geometric;
large heterophily-style boosts are not expected.
We include this experiment to stress-test Half-Hop outside the paper's
node-classification setting.

\begin{table}[H]
\centering
\caption{ModelNet10 graph classification (test accuracy \% at best validation;
seed $0$, $p=0.5$, $\alpha=0.5$).}
\label{tab:modelnet}
\begin{tabular}{lccc}
\toprule
Model & Baseline & + Half-Hop & $\Delta$ \\
\midrule
GCN & 67.62 & 68.61 & $+0.99$ \\
GraphSAGE & 69.82 & 69.27 & $-0.55$ \\
\bottomrule
\end{tabular}
\end{table}

\section{Implementation notes}
\begin{itemize}
  \item Core transform: \texttt{halfhop/halfhop.py} (aligned with
        \texttt{nerdslab/halfhop}).
  \item Supervised CLI: \texttt{experiments/supervised/run.py}.
  \item ModelNet CLI: \texttt{experiments/modelnet/run.py}
        (Kaggle dual-T4 helper: \texttt{run\_kaggle\_2gpu.sh}).
  \item For SSL Table~3 style BGRL experiments we vendor the official
        Thakoor et al.\ BGRL code and add Half-Hop as a view augmentation;
        those runs are outside the scope of the tables above.
\end{itemize}

\section{Conclusion}
We implemented Half-Hop and reproduced heterophilic supervised results and
Texas ablations in agreement with the paper's main qualitative conclusions.
As an extension, ModelNet10 shows that Half-Hop \emph{can} change accuracy on
3D mesh graphs, but the effect is modest under our single-seed protocol.
Overall, Half-Hop remains most clearly beneficial in heterophilic
node-classification settings, which matches the motivation in the original
paper.

\begin{thebibliography}{9}

\bibitem{azabou2023halfhop}
Mehdi Azabou, Venkataramana Ganesh, Shantanu Thakoor, Chi-Heng Lin,
Lakshmi Sathidevi, Ran Liu, Michal Valko, Petar Veli{\v{c}}kovi{\'c},
and Eva~L.\ Dyer.
\newblock Half-Hop: A graph upsampling approach for slowing down message
passing.
\newblock In \emph{ICML}, 2023.

\bibitem{thakoor2022bgrl}
Shantanu Thakoor, Corentin Tallec, Mohammad~Gheshlaghi Azar, Mehdi Azabou,
Eva~L.\ Dyer, R{\'e}mi Munos, Petar Veli{\v{c}}kovi{\'c}, and Michal Valko.
\newblock Large-scale representation learning on graphs via bootstrapping.
\newblock In \emph{ICLR}, 2022.

\bibitem{kipf2017gcn}
Thomas~N.\ Kipf and Max Welling.
\newblock Semi-supervised classification with graph convolutional networks.
\newblock In \emph{ICLR}, 2017.

\bibitem{hamilton2017sage}
William~L.\ Hamilton, Zhitao Ying, and Jure Leskovec.
\newblock Inductive representation learning on large graphs.
\newblock In \emph{NeurIPS}, 2017.

\bibitem{velickovic2018gat}
Petar Veli{\v{c}}kovi{\'c}, Guillem Cucurull, Arantxa Casanova, Adriana Romero,
Pietro Li{\`o}, and Yoshua Bengio.
\newblock Graph attention networks.
\newblock In \emph{ICLR}, 2018.

\bibitem{wu2015modelnet}
Zhirong Wu, Shuran Song, Aditya Khosla, Fisher Yu, Linguang Zhang,
Xiaoou Tang, and Jianxiong Xiao.
\newblock 3D ShapeNets: A deep representation for volumetric shapes.
\newblock In \emph{CVPR}, 2015.

\end{thebibliography}

\end{document}
